% TOP_PROBLEM: {"id": 3700045, "title": "Conjugate one-points in the unit disk", "category_id": 8, "category": {"id": 8, "name": "analysis", "display_name": "Analysis", "description": "Limits, continuity, calculus, and function theory.", "slug": "analysis", "order_index": 8, "created_at": "2026-07-31T15:26:25.671Z"}, "status": "partially_solved", "problem_number": "AMR-036-0045", "set_id": 13, "statusQualification": "Separates a strict two-simple-point class from its double-point limit, defines the comparison constant intrinsically, and does not add a reality condition on f."}
% FIRST_POSTED: 2026-10-01
% ENTRY_KIND: statement_only
% UnsolvedMath ID 3700045; source catalogue code AMR-036-0045
% !TeX program = lualatex
% Definition and sources only; this file is not an LLM solution attempt.
\documentclass[11pt,a4paper]{article}
\usepackage[margin=25mm]{geometry}
\usepackage{fontspec}
\setmainfont{lmroman10-regular.otf}[BoldFont=lmroman10-bold.otf,
 ItalicFont=lmroman10-italic.otf,BoldItalicFont=lmroman10-bolditalic.otf]
\usepackage{amsmath,amssymb,mathtools}
\usepackage{unicode-math}
\setmathfont{latinmodern-math.otf}
\AtBeginDocument{\let\setminus\smallsetminus}
\usepackage{xurl,hyperref}
\hypersetup{colorlinks=true,urlcolor=blue}
\setcounter{secnumdepth}{0}
\setlength{\parindent}{0pt}
\setlength{\parskip}{5pt}
\setlength{\emergencystretch}{3em}
\begin{document}
\section*{Conjugate one-points in the unit disk}\label{3700045}
{\small UnsolvedMath ID 3700045 · AMR-036-0045 · Analysis · AMR Open Problem Lists}
\subsection{Definitions and mathematical statement}
Let $\mathcal A$ consist of pairs $(f,a)$,
where $f$ is holomorphic on $\mathbb D$,
$a\in\mathbb D\setminus\mathbb R$, and
\[
 f^{-1}(0)=\{0\},\quad f'(0)\ne0,\qquad
 f^{-1}(1)=\{a,\overline a\},\quad
 f'(a)f'(\overline a)\ne0.
\]
All preimages are taken inside the disk.
Thus both $1$-points are simple and distinct.
There is no assumption that $f$ itself has
real Taylor coefficients.
Put $a_0=\inf_{(f,a)\in\mathcal A}|a|$.

For an exact comparison constant, let $c$ be
the infimum of the numbers $b\in(0,1)$
for which a holomorphic disk function has
its only zero simple at $0$ and its
only $1$-point at $b$, of multiplicity two.
This is the double-$1$-point covering extremal
constant described in the source, approximately
$0.0505468$; the approximation is not its
definition.

Determine $a_0$ and the extremal behavior.
In particular, is $a_0=c$? Decide whether
the infimum is realized with distinct
conjugate $1$-points or only approached by
a sequence in which they coalesce into a
double real $1$-point. Describe any extremizers
or limiting extremal functions.
The zero count and $1$-point count apply
globally throughout $\mathbb D$, so satisfying
the values at the displayed points alone
does not constitute an admissible construction.
\subsection{Short English statement}
Can two distinct conjugate one-points approach the origin more closely than the double-one-point covering extremal permits?
\subsection{Sources}
\begin{itemize}
\item[S1] ULAM.AI and the UnsolvedMath Contributors, supplied problems.json, record 3700045 (AMR-036-0045), AMR Open Problem Lists. Catalogue identity and source transcription; CC BY 4.0.
\url{https://huggingface.co/datasets/ulamai/UnsolvedMath}.
\item[S2] Eremenko - My favorite unsolved problems, Simultaneous stabilization, Problem 1. Original source indexed by AMR-036-0045.
\url{https://www.math.purdue.edu/~eremenko/uns1.html}.
\item[S3] A. Eremenko, Simultaneous stabilization, Problem 1. Author-maintained primary problem note.
\url{https://www.math.purdue.edu/~eremenko/dvi/simulstab.pdf}.
\end{itemize}
\end{document}
